%%======================================================================
%%  MT-1: Literature Extraction & Convention Fixing
%%        for Black Hole Entropy and Four Laws of Thermodynamics
%%
%%  This document collects formulas from the literature needed to
%%  derive the BH entropy from the SCT spectral action, compute the
%%  logarithmic correction c_log, verify all four laws, and derive
%%  modified Hawking radiation.
%%
%%  RESEARCH DOCUMENT: formulas extracted from cited papers.
%%  No original derivations --- gaps noted explicitly.
%%======================================================================
\documentclass[12pt,a4paper]{article}

%% ---- packages -------------------------------------------------------
\usepackage[utf8]{inputenc}
\usepackage[T1]{fontenc}
\usepackage{lmodern}
\usepackage{amsmath,amssymb,amsthm,mathtools}
\usepackage[margin=2.5cm]{geometry}
\usepackage{booktabs}
\usepackage{enumitem}
\usepackage{hyperref}
\usepackage{xcolor}
\usepackage{fancyhdr}
\usepackage{titlesec}
\usepackage{longtable}
\usepackage{tcolorbox}

%% ---- theorem-like environments --------------------------------------
\newtheorem{proposition}{Proposition}[section]
\newtheorem{lemma}[proposition]{Lemma}
\newtheorem{corollary}[proposition]{Corollary}
\newtheorem{theorem}[proposition]{Theorem}
\theoremstyle{definition}
\newtheorem{definition}[proposition]{Definition}
\newtheorem{convention}[proposition]{Convention}
\theoremstyle{remark}
\newtheorem{remark}[proposition]{Remark}

%% ---- custom commands ------------------------------------------------
\newcommand{\Tr}{\operatorname{Tr}}
\newcommand{\tr}{\operatorname{tr}}
\newcommand{\Rsq}{R_{\mu\nu\rho\sigma}R^{\mu\nu\rho\sigma}}
\newcommand{\Ricsq}{R_{\mu\nu}R^{\mu\nu}}
\newcommand{\Weylsq}{C_{\mu\nu\rho\sigma}C^{\mu\nu\rho\sigma}}
\newcommand{\dd}{\mathrm{d}}
\newcommand{\ii}{\mathrm{i}}
\newcommand{\Id}{\mathrm{Id}}
\newcommand{\erfi}{\operatorname{erfi}}
\DeclareMathOperator{\sgn}{sgn}

%% ---- annotation commands --------------------------------------------
\newcommand{\fromlit}[1]{\textcolor{blue!70!black}{\textsf{[#1]}}}
\newcommand{\oursign}{\textcolor{teal!80!black}{\textsf{[our convention]}}}
\newcommand{\gap}{\textcolor{red!80!black}{\textsf{[GAP]}}}
\newcommand{\needscheck}{\textcolor{orange!80!black}{\textsf{[needs check]}}}
\newcommand{\exactlit}{\textcolor{blue!60!black}{\textsf{[exact from lit.]}}}
\newcommand{\verified}[1]{%
  \par\smallskip\noindent
  \colorbox{green!10}{\parbox{\dimexpr\linewidth-2\fboxsep}{%
    \textcolor{green!50!black}{\textbf{[Verified]}} #1}}%
  \par\smallskip}
\newcommand{\signcorrection}[1]{%
  \par\noindent\colorbox{red!10}{%
    \parbox{\dimexpr\linewidth-2\fboxsep}{%
      \textcolor{red!50!black}{\textbf{[Sign correction]}}
      \textcolor{red!50!black}{#1}%
    }%
  }\par%
}

%% ---- page style -----------------------------------------------------
\pagestyle{fancy}
\fancyhf{}
\fancyhead[L]{\small\textsc{MT-1: Literature \& Conventions for
BH Entropy and Thermodynamics}}
\fancyhead[R]{\small\thepage}
\renewcommand{\headrulewidth}{0.4pt}

%% ---- hyperref -------------------------------------------------------
\hypersetup{
  colorlinks=true,
  linkcolor=blue!60!black,
  citecolor=green!40!black,
  urlcolor=blue!60!black,
  pdftitle={MT-1: Literature Extraction for BH Entropy from SCT Spectral Action},
  pdfauthor={David Alfyorov},
}

%% ---- title ----------------------------------------------------------
\title{\textbf{MT-1: Literature Extraction \& Convention Fixing\\[4pt]
       for Black Hole Entropy and Four Laws of Thermodynamics}\\[8pt]
       \large From the SCT Spectral Action to Wald Entropy,
       $c_{\log}$, and Modified Hawking Radiation}
\author{David Alfyorov}
\date{March 2026 --- Working Document}

%%======================================================================
\begin{document}
\maketitle

\begin{center}
\small\textit{Credits: Aliaksandr Samatyia contributed research-assistance
and workflow support.}
\end{center}

\thispagestyle{empty}

\begin{abstract}
This document collects all formulas, conventions, and benchmarks from the
literature needed to derive the black hole entropy from the SCT spectral
action, compute the logarithmic correction~$c_{\log}$, verify all four
laws of black hole thermodynamics, and derive the modified Hawking radiation
spectrum.  The approach proceeds via three routes: (a)~the Wald entropy
formula applied to the nonlocal spectral action (using the NT-4b variational
derivative $\delta S_{\mathrm{spec}}/\delta R_{\mu\nu\rho\sigma}$);
(b)~the Sen formula for logarithmic corrections using the SM field content
established in NT-1b Phase~3; (c)~the modified Regge--Wheeler/Zerilli
equations for greybody factors.  All formulas are extracted from the
literature with source annotations; no original derivations are performed.
Gaps requiring resolution in the derivation phase are noted explicitly.

\medskip\noindent
\textbf{Key references:}
\begin{itemize}[nosep]
\item Wald (1993): Phys.\ Rev.\ D~\textbf{48}, R3427.
  BH entropy as Noether charge.
\item Iyer--Wald (1994): Phys.\ Rev.\ D~\textbf{50}, 846.
  Properties of Noether charge and generalization to arbitrary
  diffeomorphism-invariant theories.
\item Jacobson--Myers (1993): Phys.\ Rev.\ Lett.~\textbf{70}, 3684.
  Entropy of Lovelock and higher-curvature gravity.
\item Sen (2012): arXiv:1205.0971.  Logarithmic corrections
  to Schwarzschild and other non-extremal BH entropy.
\item Wall (2015): Int.\ J.\ Mod.\ Phys.\ D~\textbf{24}, 1544014;
  also arXiv:1504.08040.  Second law for higher-derivative gravity.
\item Mann--Solodukhin (1998): Nucl.\ Phys.\ B~\textbf{523}, 293;
  arXiv:hep-th/9709064.  Conical singularity method.
\item Frolov--Fursaev (1997): Phys.\ Rev.\ D~\textbf{56}, 2212.
  Thermal fields and BH entropy.
\item Buoninfante--Mazumdar (2020): arXiv:1903.01542;
  also Buoninfante--Modesto--Shapiro~2001.07830.
  Nonlocal gravity and BH thermodynamics.
\item Connes--Chamseddine (2010): arXiv:1004.0464.
  Spectral action and gravity.
\item Stelle (1977): Phys.\ Rev.\ D~\textbf{16}, 953;
  Gen.\ Rel.\ Grav.~\textbf{9} (1978) 353.
\item Calcagni--Modesto (2022): arXiv:2211.05606.
  Nonlocal field equations.
\end{itemize}
\end{abstract}

\tableofcontents
\newpage

%%======================================================================
\section{Background: From NT-4b to MT-1}
\label{sec:background}
%%======================================================================

\subsection{What NT-4b established}

NT-4b derived the full nonlinear field equations of the SCT spectral action
on arbitrary curved backgrounds.  The key results relevant for MT-1 are:

\begin{itemize}[nosep]
  \item \textbf{Full field equations:}
    \begin{equation}\label{eq:full-eom}
      \frac{1}{\kappa^2}\,G_{\mu\nu}
      + \frac{1}{16\pi^2}\bigl[2\,F_1\!\left(\frac{\Box}{\Lambda^2}\right)
      B_{\mu\nu} + \Theta^{(C)}_{\mu\nu}\bigr]
      + \frac{1}{16\pi^2}\bigl[F_2\!\left(\frac{\Box}{\Lambda^2},\xi\right)
      H_{\mu\nu} + \Theta^{(R)}_{\mu\nu}\bigr]
      = \tfrac{1}{2}\,T_{\mu\nu}.
    \end{equation}
  \item \textbf{Ricci-flat simplification:} On Schwarzschild/Kerr,
    $R = 0$, $R_{\mu\nu} = 0$, so $H_{\mu\nu} = 0$,
    $\Theta^{(R)}_{\mu\nu} = 0$, and only the Weyl sector contributes.
  \item \textbf{Wald variation:}
    \begin{equation}\label{eq:wald-variation-nt4b}
      \frac{\partial\mathcal{L}_{\mathrm{quad}}}{\partial R_{\mu\nu\rho\sigma}}
      = \frac{1}{16\pi^2}\Bigl[
        2F_1(0)\,\frac{\partial C^2}{\partial R_{\mu\nu\rho\sigma}}
        + F_2(0,\xi)\,\frac{\partial R^2}{\partial R_{\mu\nu\rho\sigma}}
      \Bigr],
    \end{equation}
    where on the horizon $F_i(\Box/\Lambda^2) \to F_i(0)$ for
    astrophysical BH ($r_H \gg 1/\Lambda$).
  \item \textbf{Correction scale:}
    $\mathcal{O}(\Lambda^2/M_{\mathrm{Pl}}^2) \sim
    \mathcal{O}(10^{-60})$ for $\Lambda \sim \mathrm{eV}$.
\end{itemize}

\subsection{What MT-1 must achieve}

From the roadmap (Section~MT-1):
\begin{enumerate}[nosep]
  \item Derive $S_{\mathrm{BH}} = A/(4G)$ from the spectral action
    via the Wald entropy formula.
  \item Compute the logarithmic correction~$c_{\log}$ as a
    parameter-free prediction.
  \item Verify all four laws of BH thermodynamics within SCT.
  \item Derive the modified Hawking radiation spectrum (greybody factors).
\end{enumerate}

\subsection{Conventions inherited from NT-1 through NT-4b}

\begin{center}
\begin{tabular}{ll}
\toprule
Symbol & Definition \\
\midrule
$g_{\mu\nu}$ & Lorentzian metric, signature $(-,+,+,+)$ \\
$R^{\rho}{}_{\sigma\mu\nu}$ & Riemann tensor \\
$C_{\mu\nu\rho\sigma}$ & Weyl tensor \\
$\kappa^2$ & $8\pi G = 8\pi / M_{\mathrm{Pl}}^2$ \\
$\Lambda$ & Spectral cutoff scale \\
$\xi$ & Higgs non-minimal coupling to gravity \\
$\alpha_C$ & $13/120$ (total Weyl${}^2$ coefficient, parameter-free) \\
$\alpha_R(\xi)$ & $2(\xi - 1/6)^2$ (total $R^2$ coefficient) \\
$F_1(0)$ & $\alpha_C/(16\pi^2) = 13/(1920\pi^2)$ \\
$F_2(0,\xi)$ & $\alpha_R(\xi)/(16\pi^2)$ \\
$m_2$ & $\Lambda\sqrt{60/13}$ (spin-2 effective mass) \\
$m_0(\xi)$ & $\Lambda/\sqrt{6(\xi - 1/6)^2}$ (spin-0 effective mass) \\
$\varphi(x)$ & $e^{-x/4}\sqrt{\pi/x}\,\erfi(\sqrt{x}/2)$
  (master function, entire) \\
\bottomrule
\end{tabular}
\end{center}

\begin{convention}[SM field content]\label{conv:SM}
\fromlit{CPR 0805.2909; NT-1b Phase~3.}
\begin{center}
\begin{tabular}{lccc}
\toprule
Field & Spin & Count & Convention \\
\midrule
Real scalars (Higgs) & 0 & $N_s = 4$ & Real components \\
Weyl fermions & 1/2 & $N_f = 45$ & Weyl count \\
Dirac fermions & 1/2 & $N_D = N_f/2 = 22.5$ & Dirac count \\
Gauge bosons & 1 & $N_v = 12$ & $SU(3) \times SU(2) \times U(1)$ \\
\bottomrule
\end{tabular}
\end{center}
\end{convention}


%%======================================================================
\section{Schwarzschild Geometry}
\label{sec:schwarzschild}
%%======================================================================

\subsection{Metric and curvature}
\fromlit{MTW (1973); Wald (1984), Chapter~6.}

The Schwarzschild metric in standard coordinates:
\begin{equation}\label{eq:schw-metric}
  \dd s^2 = -f(r)\,\dd t^2 + f(r)^{-1}\,\dd r^2
  + r^2\,\dd\Omega^2,
  \qquad f(r) = 1 - \frac{r_H}{r},
  \qquad r_H = 2GM.
\end{equation}

\begin{proposition}[Schwarzschild curvature invariants]
\label{prop:schw-curv}
\fromlit{Wald (1984), eq.~(6.1.7)--(6.1.11).}
\begin{align}
  R &= 0, \qquad R_{\mu\nu} = 0 \quad (\text{vacuum}), \\
  R_{\mu\nu\rho\sigma}R^{\mu\nu\rho\sigma}
  &= \frac{48\,G^2 M^2}{r^6}, \label{eq:kretschner}\\
  C_{\mu\nu\rho\sigma} &= R_{\mu\nu\rho\sigma}
  \quad (\text{since Ricci-flat}), \\
  \Weylsq \Big|_{r=r_H}
  &= \frac{48\,G^2 M^2}{(2GM)^6} = \frac{3}{4\,G^4 M^4}.
  \label{eq:weyl-sq-horizon}
\end{align}
\end{proposition}

\subsection{Horizon properties}
\fromlit{Wald (1984); Poisson (2004).}

\begin{center}
\begin{tabular}{ll}
\toprule
Quantity & Value \\
\midrule
Horizon radius & $r_H = 2GM$ \\
Surface gravity & $\kappa = f'(r_H)/2 = 1/(4GM)$ \\
Area & $A = 4\pi r_H^2 = 16\pi G^2 M^2$ \\
Hawking temperature & $T_H = \kappa/(2\pi) = 1/(8\pi GM)$ \\
Bekenstein--Hawking entropy & $S_{\mathrm{BH}} = A/(4G) = 4\pi G M^2$ \\
\bottomrule
\end{tabular}
\end{center}

\subsection{Bifurcation surface and binormal}
\fromlit{Wald (1993), eq.~(12)--(14); Iyer--Wald (1994), eq.~(8).}

The bifurcation 2-sphere $\mathcal{H}$ is the $(d-2)$-surface at which the
Killing horizon of $\xi^a = (\partial/\partial t)^a$ has vanishing norm.
The binormal $\hat{\epsilon}_{\mu\nu}$ is the antisymmetric tensor normal to
$\mathcal{H}$, lying in the $(t,r)$ plane:
\begin{equation}\label{eq:binormal}
  \hat{\epsilon}_{\mu\nu} = 2\,n_{[\mu}\,s_{\nu]},
  \qquad \hat{\epsilon}_{\mu\nu}\hat{\epsilon}^{\mu\nu} = -2,
\end{equation}
where $n^\mu$ is the future-directed unit timelike normal and $s^\mu$ is
the outward-directed unit spacelike normal to~$\mathcal{H}$.

\begin{proposition}[Weyl contraction on Schwarzschild bifurcation surface]
\label{prop:weyl-binormal}
\fromlit{Jacobson--Myers (1993); Padmanabhan (2010), eq.~(15.62).}

In the orthonormal frame $(e^{\hat{0}}, e^{\hat{1}},
e^{\hat{2}}, e^{\hat{3}})$ adapted to the $(t,r)$ binormal:
\begin{equation}\label{eq:C-binormal}
  C_{\mu\nu\rho\sigma}\,\hat{\epsilon}^{\mu\nu}\hat{\epsilon}^{\rho\sigma}
  \Big|_{\mathcal{H}}
  = 4\,C_{\hat{0}\hat{1}\hat{0}\hat{1}}\Big|_{r=r_H}
  = 4 \cdot \Bigl(-\frac{2GM}{r_H^3}\Bigr)
  = -\frac{1}{G^2 M^2}.
\end{equation}
\end{proposition}

\begin{remark}[Uniformity on the 2-sphere]
The Schwarzschild Weyl tensor has $\ell = 0$ angular dependence in the
orthonormal frame (it is spherically symmetric), so the integrand in the
Wald formula is constant over $\mathcal{H}$.  The integral over $S^2$
simply multiplies by the area $A = 4\pi r_H^2$.
\end{remark}


%%======================================================================
\section{The Wald Entropy Formula}
\label{sec:wald}
%%======================================================================

\subsection{General statement}
\fromlit{Wald (1993), Phys.\ Rev.\ D~\textbf{48}, R3427;
Iyer--Wald (1994), Phys.\ Rev.\ D~\textbf{50}, 846.}

For any diffeomorphism-invariant Lagrangian $L(g_{\mu\nu},
R_{\mu\nu\rho\sigma}, \nabla_\alpha R_{\mu\nu\rho\sigma}, \ldots)$,
the black hole entropy is
\begin{equation}\label{eq:wald-general}
  \boxed{
  S_{\mathrm{Wald}} = -2\pi \oint_{\mathcal{H}}
  \frac{\partial\mathcal{L}}{\partial R_{\mu\nu\rho\sigma}}\,
  \hat{\epsilon}_{\mu\nu}\,\hat{\epsilon}_{\rho\sigma}\,\dd A.}
\end{equation}
This formula gives the correct first law of black hole mechanics for
\emph{all} diffeomorphism-invariant theories.

\begin{remark}[Covariant derivatives of Riemann]
When the Lagrangian depends on $\nabla_\alpha R_{\mu\nu\rho\sigma}$
(as in our nonlocal spectral action), the Wald formula generalizes
via integration by parts.  For a Lagrangian of the form
$\mathcal{L} = f(\Box)\,I$ where $I$ is a curvature invariant and
$\Box = g^{\mu\nu}\nabla_\mu\nabla_\nu$, the on-horizon evaluation
uses the local limit $f(\Box) \to f(0)$ for astrophysical BH
($r_H \gg 1/\Lambda$), and additional $\nabla$-dependent corrections
are suppressed by powers of $\kappa^2/\Lambda^2 \ll 1$.
\fromlit{Buoninfante--Mazumdar (2020); Conroy--Edholm--Mazumdar (2017),
arXiv:1509.01636.}
\end{remark}

\subsection{Einstein--Hilbert contribution}
\fromlit{Wald (1993), eq.~(14); Iyer--Wald (1994), eq.~(16).}

For $\mathcal{L}_{\mathrm{EH}} = R/(16\pi G)$:
\begin{equation}
  \frac{\partial\mathcal{L}_{\mathrm{EH}}}
  {\partial R_{\mu\nu\rho\sigma}}
  = \frac{1}{16\pi G}\,
  \frac{\partial R}{\partial R_{\mu\nu\rho\sigma}}
  = \frac{1}{16\pi G}\,
  g^{\mu[\rho}\,g^{\sigma]\nu}.
\end{equation}

Contracting with the binormal:
\begin{equation}
  g^{\mu[\rho}\,g^{\sigma]\nu}\,
  \hat{\epsilon}_{\mu\nu}\,\hat{\epsilon}_{\rho\sigma}
  = -2,
\end{equation}
so
\begin{equation}
  S_{\mathrm{EH}} = -2\pi \cdot \frac{1}{16\pi G} \cdot (-2)
  \cdot A = \frac{A}{4G}.
\end{equation}
This is the standard Bekenstein--Hawking result.

\subsection{Weyl-squared ($C^2$) contribution}
\label{sec:wald-weyl}
\fromlit{Jacobson--Myers (1993), eq.~(7);
Jacobson--Kang--Myers (1994), arXiv:gr-qc/9312023.}

For $\mathcal{L}_{C^2} = \alpha\,\Weylsq$ where
$\alpha = \alpha_C/(16\pi^2)$ \oursign:
\begin{equation}
  \frac{\partial\mathcal{L}_{C^2}}
  {\partial R_{\mu\nu\rho\sigma}}
  = \alpha\,\frac{\partial(\Weylsq)}
  {\partial R_{\mu\nu\rho\sigma}}.
\end{equation}

The Weyl tensor decomposition is
\begin{equation}\label{eq:weyl-decomp}
  C_{\mu\nu\rho\sigma} = R_{\mu\nu\rho\sigma}
  - \bigl(g_{\mu[\rho}\,R_{\sigma]\nu}
  - g_{\nu[\rho}\,R_{\sigma]\mu}\bigr)
  + \tfrac{1}{3}\,R\,g_{\mu[\rho}\,g_{\sigma]\nu}.
\end{equation}

On a Ricci-flat background ($R = 0$, $R_{\mu\nu} = 0$):
\begin{equation}
  C_{\mu\nu\rho\sigma}\big|_{\mathrm{Ricci\text{-}flat}}
  = R_{\mu\nu\rho\sigma},
  \qquad
  \Weylsq\big|_{\mathrm{Ricci\text{-}flat}}
  = \Rsq.
\end{equation}

\begin{proposition}[Jacobson--Myers formula on Ricci-flat backgrounds]
\label{prop:JM}
\fromlit{Jacobson--Myers (1993), eq.~(7)--(10).}

For $\mathcal{L} = \gamma\,\Rsq$ (Gauss-Bonnet-type term
in the Riemann-squared sector), the Wald entropy correction is
\begin{equation}\label{eq:JM-formula}
  \delta S_{\mathrm{JM}} = 8\pi\gamma\,\chi(\mathcal{H}),
\end{equation}
where $\chi(\mathcal{H})$ is the Euler characteristic of the horizon
cross-section.  For $\mathcal{H} = S^2$: $\chi = 2$.
\end{proposition}

\begin{corollary}[Wald entropy correction from $C^2$ on Schwarzschild]
\label{cor:delta-S-C2}
On Schwarzschild, $\Weylsq = \Rsq$ and the coefficient of
$\Rsq$ in the action is $\gamma = \alpha = \alpha_C/(16\pi^2)$.
The Jacobson--Myers formula~\eqref{eq:JM-formula} gives:
\begin{equation}\label{eq:delta-S-C2}
  \boxed{
  \delta S_{C^2} = 8\pi \cdot \frac{\alpha_C}{16\pi^2} \cdot 2
  = \frac{\alpha_C}{\pi}
  = \frac{13}{120\pi}
  \approx 0.0345.}
\end{equation}
This is a \textbf{topological} correction (independent of the horizon
area~$A$), arising from the Euler characteristic of~$S^2$.
\end{corollary}

\begin{remark}[Area-independent correction]
The $C^2$ Wald entropy correction is $\alpha_C/\pi \approx 0.035$,
which is a dimensionless constant of order unity.  For an astrophysical
black hole with $S_{\mathrm{BH}} \sim 10^{79}$, this correction is
utterly negligible.  The correction becomes relevant only for
micro-black holes with $A \sim \ell_P^2$, where higher-order terms
also become important.
\end{remark}

\subsection{$R^2$ contribution}
\fromlit{Jacobson--Myers (1993); Padmanabhan (2010), Chapter~15.}

For $\mathcal{L}_{R^2} = \beta\,R^2$ with
$\beta = \alpha_R(\xi)/(16\pi^2)$:
\begin{equation}
  \frac{\partial(R^2)}{\partial R_{\mu\nu\rho\sigma}}
  = 2R\,g^{\mu[\rho}\,g^{\sigma]\nu}.
\end{equation}

\begin{proposition}[$R^2$ contribution vanishes on Schwarzschild]
\label{prop:R2-vanishes}
On Schwarzschild, $R = 0$, so
\begin{equation}
  \delta S_{R^2}\big|_{\mathrm{Schw}} = 0.
\end{equation}
The $R^2$ sector contributes to the Wald entropy only on backgrounds
with nonzero Ricci scalar (de~Sitter, Kerr--de~Sitter,
Reissner--Nordstr\"om--de~Sitter, etc.).
\end{proposition}

\subsection{Nonlocal corrections}
\label{sec:nonlocal-wald}
\gap

The full SCT action contains nonlocal form factors
$F_i(\Box/\Lambda^2)$, and the Wald formula for nonlocal actions
requires additional care.

\begin{proposition}[Local limit on astrophysical horizons]
\label{prop:local-limit}
\fromlit{Conroy--Edholm--Mazumdar (2017), arXiv:1509.01636, Section~5;
Buoninfante--Mazumdar (2020), arXiv:1903.01542, Section~3.}

For a Schwarzschild BH with mass $M \gg M_{\mathrm{Pl}}$:
\begin{itemize}[nosep]
  \item The curvature scale on the horizon is
    $\kappa_H \sim 1/(GM) \ll \Lambda$.
  \item The d'Alembertian eigenvalue on the horizon satisfies
    $\Box/\Lambda^2 \sim \kappa_H^2/\Lambda^2 \ll 1$.
  \item The form factors reduce to their local ($z = 0$) values:
    $F_i(\Box/\Lambda^2) \to F_i(0)$.
  \item Nonlocal corrections to the Wald entropy are suppressed
    by $(\kappa_H/\Lambda)^2 \sim (M_{\mathrm{Pl}}/M)^2
    \cdot (\Lambda/M_{\mathrm{Pl}})^2$.
\end{itemize}
\end{proposition}

\begin{remark}[Near-Planckian regime]
For micro-BH with $M \sim M_{\mathrm{Pl}}$ or
$r_H \sim 1/\Lambda$, the nonlocal structure of $F_1(z)$
becomes important and the local Wald formula does not apply directly.
In this regime, the entire-function property of $F_1(z)$ (from NT-2)
prevents pole-induced divergences, but the full nonlocal Wald formula
must be used.  \gap\ \textit{The full nonlocal Wald formula for
$\int F(\Box)\,C^2$ requires resolution in the derivation phase.}
\end{remark}


%%======================================================================
\section{Full Wald Entropy in SCT}
\label{sec:full-wald}
%%======================================================================

Combining the Einstein--Hilbert contribution (\S\ref{sec:wald}),
the $C^2$ correction (Corollary~\ref{cor:delta-S-C2}), and the
vanishing $R^2$ contribution (Proposition~\ref{prop:R2-vanishes}):

\begin{theorem}[Wald entropy of Schwarzschild BH in SCT]
\label{thm:wald-schw}
For a Schwarzschild BH in the SCT spectral action, the Wald entropy is
\begin{equation}\label{eq:S-wald-full}
  \boxed{
  S_{\mathrm{Wald}} = \frac{A}{4G}
  + \frac{\alpha_C}{\pi}
  = \frac{A}{4G} + \frac{13}{120\pi}.}
\end{equation}
The first term is the Bekenstein--Hawking area law.  The second term is
a topological correction from the Weyl-squared sector,
determined entirely by the SM field content via $\alpha_C = 13/120$.
\end{theorem}

For Kerr--Newman BH with nonzero charge/angular momentum:
\begin{equation}
  S_{\mathrm{Wald}} = \frac{A}{4G}
  + \frac{\alpha_C}{\pi}\,\chi(\mathcal{H})
  + \frac{\alpha_R(\xi)}{\pi}\,\delta S_{R^2}(\xi),
\end{equation}
where $\chi(\mathcal{H}) = 2$ for topologically spherical horizons and
$\delta S_{R^2}$ is nonzero only when $R \neq 0$
(e.g., in de~Sitter backgrounds).


%%======================================================================
\section{Logarithmic Correction $c_{\log}$}
\label{sec:clog}
%%======================================================================

\subsection{Sen's formula}
\fromlit{Sen (2012), arXiv:1205.0971, eq.~(1.2) and Table~1.}

For a 4-dimensional Schwarzschild BH in Einstein gravity coupled to
$n_S$~real scalars, $n_F$~massless \textbf{Dirac} fermions, and
$n_V$~massless vectors, the one-loop logarithmic correction to the
microcanonical entropy is
\begin{equation}\label{eq:sen-formula}
  \boxed{
  c_{\log} = \frac{1}{180}\bigl[
    2\,n_S + 7\,n_F - 26\,n_V + 424
  \bigr].}
\end{equation}
The $+424$ accounts for the graviton contribution (including zero modes).
The vector coefficient is $-26$ (negative, due to Faddeev--Popov ghost
dominance in the Kretschner term of the Seeley--DeWitt $a_4$ coefficient).

\signcorrection{Corrected from LR audit.  The original extraction had
the vector and graviton signs reversed and used Weyl fermion counting
with coefficient~7.  The correct Sen formula uses \textbf{Dirac}
fermions with coefficient~7, vector coefficient $-26$, and graviton
$+424$.  Verified against Sen (2012) eq.~(1.2) and independently
against arXiv:2104.14902 eq.~(5).}

\begin{remark}[Derivation method]
Sen's result is obtained via the Euclidean gravity method (heat kernel
on the Euclidean Schwarzschild instanton), with careful treatment of
the zero modes of the graviton kinetic operator.  The computation
uses the Seeley--DeWitt coefficients $a_4(x)$ evaluated on the
instanton background, integrated against the logarithmic divergence
$\ln(\Lambda^2 A)$.  The individual field contributions are:
\begin{center}
\begin{tabular}{lc}
\toprule
Field & Contribution to $180\,c_{\log}$ \\
\midrule
Real scalar & $+2$ per field \\
Dirac fermion & $+7$ per field \\
Massless vector & $-26$ per field \\
Graviton & $+424$ (total) \\
\bottomrule
\end{tabular}
\end{center}
\end{remark}

\begin{remark}[Weyl fermion convention]
The formula can be rewritten in terms of Weyl fermions
($n_W = 2\,n_F$):
\begin{equation}
  c_{\log} = \frac{1}{180}\bigl[
    2\,n_S + \tfrac{7}{2}\,n_W - 26\,n_V + 424
  \bigr].
\end{equation}
This is equivalent to~\eqref{eq:sen-formula} by the substitution
$n_F = n_W/2$.  \textbf{Note:} the coefficient per Weyl fermion
is $7/2$, not~$7$.
\end{remark}

\subsection{SCT prediction for $c_{\log}$}
\label{sec:clog-sct}

Using the SM field content from Convention~\ref{conv:SM}
($n_S = 4$, $n_F = N_D = 45/2 = 22.5$ Dirac fermions, $n_V = 12$):

\begin{theorem}[Parameter-free prediction of $c_{\log}$ in SCT]
\label{thm:clog}
\begin{equation}\label{eq:clog-sct}
  \boxed{
  c_{\log} = \frac{1}{180}\bigl[
    2 \cdot 4 + 7 \cdot \tfrac{45}{2} - 26 \cdot 12 + 424
  \bigr]
  = \frac{37}{24}
  \approx 1.542.}
\end{equation}
This is the \textbf{microcanonical} coefficient (no zero-mode
subtraction).
\end{theorem}

\begin{proof}[Verification]
$2 \cdot 4 = 8$; $7 \cdot 22.5 = 157.5$; $-26 \cdot 12 = -312$;
graviton: $+424$.  Total $= 8 + 157.5 - 312 + 424 = 277.5$.
$c_{\log} = 277.5/180 = 37/24$.
\end{proof}

\begin{remark}[Matter vs.\ graviton decomposition]
The graviton contributes $+424/180 = 106/45 \approx 2.356$
(positive: gravity \emph{increases} the logarithmic correction).
The matter contribution is $(8 + 157.5 - 312)/180 = -146.5/180
\approx -0.814$ (net negative because the 12~vectors dominate
over scalars and fermions).
The net result is positive: $c_{\log} = 37/24 > 0$, meaning quantum
corrections \emph{increase} the entropy relative to the area law.
\end{remark}

\begin{remark}[Comparison with competing approaches]
\begin{center}
\begin{tabular}{lcc}
\toprule
Approach & $c_{\log}$ (4D Schwarzschild) & Reference \\
\midrule
\textbf{SCT (this work)} & $\mathbf{37/24 \approx 1.542}$
  & \textbf{parameter-free} \\
LQG (isolated horizons) & $-3/2 = -1.5$
  & Kaul--Majumdar (2000) \\
String theory & model-dependent
  & Sen (2012) \\
Asymptotic safety & not well-defined
  & (running $G$) \\
IDG (Modesto) & same as Sen formula
  & Modesto (2012) \\
\bottomrule
\end{tabular}
\end{center}

\signcorrection{The LQG prediction $c_{\log} = -3/2$ has the
\emph{opposite sign} from the SCT prediction.  This is a potentially
discriminating observable if logarithmic corrections to BH entropy
become accessible, though this is far beyond current experimental
capability.  Sen (2012) notes that the LQG result disagrees with the
Euclidean gravity computation.}
\end{remark}

\subsection{Relation to conformal anomaly}
\fromlit{Christensen--Duff (1978), Nucl.\ Phys.\ B~\textbf{154}, 301.}

The coefficients in Sen's formula are related to the
Seeley--DeWitt coefficient $a_4$, which in turn determines the
conformal (trace) anomaly:
\begin{equation}
  \langle T^\mu{}_\mu \rangle
  = \frac{1}{16\pi^2}\bigl[c\,\Weylsq - a\,E_4
  + b\,\Box R\bigr],
\end{equation}
where $E_4 = \Rsq - 4\Ricsq + R^2$ is the Euler density and
$a$, $c$ are the central charges.  The relationship to our
$\alpha_C$ is:
\begin{equation}
  c = \frac{1}{120}\bigl[N_s + 6N_D + 12N_v\bigr],
  \qquad
  a = \frac{1}{360}\bigl[N_s + \tfrac{11}{2}\cdot 2N_D + 62N_v\bigr].
\end{equation}

\needscheck\ The exact mapping between $c_{\log}$ and $(a, c)$ involves
additional contributions from the graviton sector and zero-mode
subtleties.  The Sen formula~\eqref{eq:sen-formula} incorporates all
these corrections; it does not simply reduce to $c_{\log} = f(a, c)$.


%%======================================================================
\section{Four Laws of BH Thermodynamics in SCT}
\label{sec:four-laws}
%%======================================================================

\subsection{Zeroth law: constancy of surface gravity}
\fromlit{Bardeen--Carter--Hawking (1973), Commun.\ Math.\ Phys.\
\textbf{31}, 161; Wald (1984), Section~12.5.}

\begin{theorem}[Zeroth law -- standard GR]
For a stationary black hole satisfying the dominant energy condition,
the surface gravity $\kappa$ is constant on the event horizon.
\end{theorem}

\textbf{Status in SCT:} The zeroth law depends on the field equations
through the dominant energy condition (DEC).  The SCT field
equations~\eqref{eq:full-eom} modify the Einstein equations, but
for Ricci-flat backgrounds (Schwarzschild, Kerr), the modifications
are $\mathcal{O}(\Lambda^2/M_{\mathrm{Pl}}^2)$, and the zeroth
law holds to the same accuracy.

\gap\ For non-vacuum BH (e.g., charged or hairy), verification
requires checking the DEC for the effective stress-energy tensor
including the higher-derivative contributions.

\subsection{First law: energy balance}
\fromlit{Wald (1993); Iyer--Wald (1994), eq.~(27)--(30).}

For the SCT action, the first law takes the form
\begin{equation}\label{eq:first-law}
  \dd M = \frac{\kappa_{\mathrm{eff}}}{2\pi}\,
  \dd S_{\mathrm{Wald}} + \Omega_H\,\dd J + \Phi_H\,\dd Q,
\end{equation}
where $S_{\mathrm{Wald}}$ is the full Wald entropy~\eqref{eq:S-wald-full}
and $\kappa_{\mathrm{eff}}$ is the effective surface gravity.

\begin{proposition}[First law for Wald entropy]
\fromlit{Iyer--Wald (1994), Theorem~1.}
For any diffeomorphism-invariant Lagrangian, the first law
$\dd M = (\kappa/2\pi)\,\dd S_{\mathrm{Wald}} + \cdots$ holds
\emph{exactly}, provided one uses the Wald entropy (not the area).
This is a consequence of the Noether charge construction and does
not require the DEC.
\end{proposition}

\begin{remark}[Modified first law in NCG]
A 2025 result (arXiv:2509.23423, referenced in the SCT roadmap)
shows that the conventional first law is violated in noncommutative
geometry because the entropy deviates from Bekenstein--Hawking.
In SCT, the first law is \emph{maintained} precisely because we
use the Wald entropy, which is the Noether charge.  The ``violation''
in NCG arises from using the area entropy $A/(4G)$ instead of the
Wald entropy.
\end{remark}

\subsection{Second law: area non-decrease}
\fromlit{Wall (2015), arXiv:1504.08040;
Wall (2024), arXiv:2312.17636.}

\begin{theorem}[Wall's second law for higher-derivative gravity]
\fromlit{Wall (2015), Theorem~1; Wall (2024).}
For a diffeomorphism-invariant theory of gravity in $d \geq 4$
dimensions, the generalized second law
\begin{equation}
  \dd S_{\mathrm{gen}} = \dd S_{\mathrm{Wald}}
  + \dd S_{\mathrm{outside}} \geq 0
\end{equation}
holds provided:
\begin{enumerate}[nosep]
  \item The theory has a well-posed initial value problem.
  \item The null energy condition holds for the matter sector.
  \item The entropy functional satisfies certain linearized stability
    conditions.
\end{enumerate}
\end{theorem}

\textbf{Status in SCT:}
\begin{enumerate}[nosep]
  \item \textbf{Well-posedness:} The SCT field equations include
    higher-derivative terms, which generically threaten well-posedness.
    However, the entire-function property of $F_1(z)$
    (NT-2, ghost-freedom) and the MR-3 causality analysis suggest
    that the initial value problem is well-posed at least for
    macroscopic perturbations.  \gap\ Full proof of well-posedness
    for the nonlocal SCT equations is an open problem (related to MR-4).
  \item \textbf{Null energy condition:} The SM matter sector satisfies
    the NEC classically.
  \item \textbf{Linearized stability:} The spin-2 propagator
    denominator $\Pi_{\mathrm{TT}}(z)$ has ghost poles (MR-2),
    which may violate the stability condition.
    \gap\ This requires careful analysis in the derivation phase.
\end{enumerate}

\signcorrection{The second law is CONDITIONAL on the ghost resolution
(MR-2).  If the ghosts are resolved via the Lee--Wick/fakeon
prescription (FK analysis), the effective theory may satisfy Wall's
conditions in the macroscopic regime where ghost effects are
suppressed by $(M_{\mathrm{Pl}}/\Lambda)^2$.}

\subsection{Third law: extremal limit}
\fromlit{Israel (1986), Phys.\ Rev.\ Lett.~\textbf{57}, 397;
Wald (1997), arXiv:gr-qc/9702022.}

\begin{theorem}[Third law -- Nernst version]
It is impossible to reduce the surface gravity $\kappa$ of a black
hole to zero by any finite physical process.
\end{theorem}

\textbf{Status in SCT:} For the Schwarzschild BH, $\kappa = 1/(4GM)$
vanishes only as $M \to \infty$, which is not an extremal limit.
For Kerr--Newman BH, the extremal condition $\kappa = 0$ occurs at
$M^2 = a^2 + Q^2$ (where $a = J/M$).  The SCT corrections to the
extremal condition are:
\begin{equation}
  M_{\mathrm{ext}}^2 = a^2 + Q^2
  + \mathcal{O}\!\left(\frac{\Lambda^2}{M_{\mathrm{Pl}}^2}\right)
  \cdot (a^2 + Q^2).
\end{equation}
The correction is negligible for astrophysical parameters.

\gap\ The Wald entropy at extremality and the attractor mechanism
in the SCT context require further analysis (connects to MR-9).


%%======================================================================
\section{Modified Hawking Radiation}
\label{sec:hawking}
%%======================================================================

\subsection{Standard Hawking spectrum}
\fromlit{Hawking (1975), Commun.\ Math.\ Phys.~\textbf{43}, 199;
Page (1976), Phys.\ Rev.\ D~\textbf{13}, 198.}

The Hawking temperature of a Schwarzschild BH is
\begin{equation}
  T_H = \frac{\kappa}{2\pi} = \frac{1}{8\pi GM}.
\end{equation}

The emission rate per unit frequency for a species of spin~$s$ is
\begin{equation}\label{eq:hawking-rate}
  \frac{\dd N^{(s)}}{\dd\omega\,\dd t}
  = \frac{1}{2\pi}\,\frac{\Gamma^{(s)}_\ell(\omega)}
  {e^{\omega/T_H} \pm 1},
\end{equation}
where $\Gamma^{(s)}_\ell(\omega)$ is the greybody factor (absorption
probability) for partial wave~$\ell$ and spin~$s$, and $+$ applies
to fermions, $-$ to bosons.

\subsection{Greybody factors and the Regge--Wheeler equation}
\fromlit{Regge--Wheeler (1957), Phys.\ Rev.~\textbf{108}, 1063;
Zerilli (1970), Phys.\ Rev.\ D~\textbf{2}, 2141.}

In the standard Regge--Wheeler formalism, the radial wave equation
for spin-$s$ perturbations of Schwarzschild is
\begin{equation}\label{eq:RW}
  \frac{\dd^2 \Psi}{\dd r_*^2}
  + \bigl[\omega^2 - V^{(s)}_{\mathrm{eff}}(r)\bigr]\Psi = 0,
\end{equation}
where $r_* = r + r_H\ln(r/r_H - 1)$ is the tortoise coordinate and
\begin{equation}
  V^{(s)}_{\mathrm{eff}}(r) = f(r)\left[
    \frac{\ell(\ell+1)}{r^2}
    + \frac{(1-s^2)\,r_H}{r^3}
  \right].
\end{equation}

\subsection{SCT modifications to the effective potential}
\gap

In SCT, the linearized equations around the Schwarzschild background
are modified by the nonlocal form factors.  The spin-2 perturbation
equation becomes (schematically):
\begin{equation}\label{eq:RW-SCT}
  \Pi_{\mathrm{TT}}\!\left(\frac{-\partial_t^2 + \nabla_r^2}
  {\Lambda^2}\right)\!
  \left[\frac{\dd^2 \Psi}{\dd r_*^2}
  + \omega^2\,\Psi\right]
  - V^{(2)}_{\mathrm{eff}}(r)\,\Psi = 0.
\end{equation}

\begin{proposition}[SCT modifications to greybody factors]
\label{prop:greybody}
\fromlit{Conroy--Koshelev--Mazumdar (2017), arXiv:1605.02080, Section~4;
Buoninfante et al.\ (2020), arXiv:1906.11238, Section~3.}
\begin{itemize}[nosep]
  \item At $\omega \ll \Lambda$: $\Pi_{\mathrm{TT}} \to 1$ and the
    standard Regge--Wheeler equation is recovered.  Modifications to
    greybody factors are exponentially suppressed as
    $\sim e^{-\Lambda r_H}$.
  \item At $\omega \sim \Lambda$: the nonlocal form factor modifies
    the effective potential, potentially suppressing high-energy
    emission.
  \item For astrophysical BH ($M \gg M_{\mathrm{Pl}}$):
    $T_H \sim 10^{-8}\,\mathrm{K} \cdot (M_\odot/M) \ll \Lambda$,
    so all emitted quanta satisfy $\omega \ll \Lambda$ and SCT
    corrections are completely negligible.
  \item The modified evaporation endpoint (where $T_H \sim \Lambda$
    at $M \sim M_{\mathrm{Pl}}^2/\Lambda$) is where SCT corrections
    become significant.  This connects to MR-9 (singularity
    resolution and information paradox).
\end{itemize}
\end{proposition}


%%======================================================================
\section{Entanglement Entropy Approach}
\label{sec:entanglement}
%%======================================================================

\fromlit{Bombelli et al.\ (1986), Phys.\ Rev.\ D~\textbf{34}, 373;
Srednicki (1993), Phys.\ Rev.\ Lett.~\textbf{71}, 666.}

An alternative route to the BH entropy uses the entanglement entropy
of quantum fields across the horizon:
\begin{equation}
  S_{\mathrm{ent}} = -\Tr(\rho_{\mathrm{in}}\ln\rho_{\mathrm{in}}),
  \qquad
  \rho_{\mathrm{in}} = \Tr_{\mathrm{out}}
  |\Omega\rangle\langle\Omega|,
\end{equation}
where $|\Omega\rangle$ is the vacuum state and the trace is over
degrees of freedom outside the horizon.

\begin{proposition}[Entanglement entropy and area law]
\fromlit{Bombelli et al.\ (1986); Srednicki (1993).}
For a free scalar field on a flat background with a spherical
entangling surface:
\begin{equation}
  S_{\mathrm{ent}} = c_1\,\frac{A}{\epsilon^2}
  + c_2\,\ln\!\left(\frac{A}{\epsilon^2}\right)
  + \mathcal{O}(1),
\end{equation}
where $\epsilon$ is the UV cutoff.  The area-law term is UV-divergent;
the logarithmic term has a universal (cutoff-independent) coefficient.
\end{proposition}

\textbf{Connection to SCT:} In the spectral action framework, the
UV cutoff is naturally provided by the spectral scale~$\Lambda$.
The entanglement entropy of the Dirac operator across the horizon
should reproduce the Wald entropy in the semiclassical limit.

\gap\ The explicit computation of $S_{\mathrm{ent}}$ for the SCT
spectral triple restricted to a causal diamond requires the
near-horizon spectral geometry, which is an output of MT-1 (not
an input).  This is the ``spectral action on horizon'' approach
(b) in the roadmap.


%%======================================================================
\section{Benchmarks and Cross-Checks}
\label{sec:benchmarks}
%%======================================================================

\begin{center}
\begin{longtable}{lllc}
\toprule
\textbf{Check} & \textbf{Expected} & \textbf{Source} & \textbf{Status} \\
\midrule
\endhead
$S_{\mathrm{EH}} = A/(4G)$ & Bekenstein--Hawking & Wald (1993) &
  \textcolor{green!50!black}{PASS} \\
$\delta S_{C^2} = \alpha_C/\pi$ & $13/(120\pi) \approx 0.035$ &
  Jacobson--Myers (1993) & \textcolor{orange}{To verify} \\
$\delta S_{R^2}\big|_{\mathrm{Schw}} = 0$ & $R = 0$ &
  Direct & \textcolor{green!50!black}{PASS} \\
$c_{\log} = 37/24$ & $\approx 1.542$ & Sen (2012) &
  \textcolor{orange}{To verify} \\
First law holds with $S_{\mathrm{Wald}}$ & Iyer--Wald (1994) &
  & \textcolor{orange}{To verify} \\
Zeroth law: $\kappa = \mathrm{const}$ & Standard GR proof &
  BCH (1973) & \textcolor{orange}{Conditional} \\
Greybody $\to$ standard at $\omega \ll \Lambda$ & Regge--Wheeler &
  & \textcolor{orange}{To verify} \\
\bottomrule
\end{longtable}
\end{center}


%%======================================================================
\section{Gaps and Open Questions}
\label{sec:gaps}
%%======================================================================

\begin{enumerate}[label=\textbf{G\arabic*:},leftmargin=3em]
  \item \textbf{Nonlocal Wald formula.}
    The full Wald formula for a nonlocal Lagrangian
    $\int F(\Box)\,C^2$ requires integration by parts of the
    form factor away from the local limit.  For astrophysical BH
    the local limit suffices, but for micro-BH this is essential.
    \gap

  \item \textbf{Weyl binormal contraction.}
    The contraction $C_{\mu\nu\rho\sigma}\,\hat{\epsilon}^{\mu\nu}
    \hat{\epsilon}^{\rho\sigma}$ on the Schwarzschild bifurcation
    surface must be computed explicitly in a regular coordinate
    system (Kruskal) and verified against the orthonormal-frame
    result~\eqref{eq:C-binormal}.  \needscheck

  \item \textbf{Second law and ghost sector.}
    Wall's proof requires well-posedness and linearized stability.
    The SCT ghost poles (MR-2) may violate these conditions.
    Status: CONDITIONAL on ghost resolution.  \gap

  \item \textbf{Entanglement entropy.}
    The spectral triple approach to entanglement entropy across
    the horizon has not been computed explicitly for SCT.  This
    is the second route (b) in the roadmap.  \gap

  \item \textbf{Kerr generalization.}
    The Kerr BH has nonzero frame-dragging, and the Weyl binormal
    contraction differs from Schwarzschild.  The $R^2$ sector still
    vanishes (Kerr is vacuum), but the topological correction
    $\alpha_C/\pi$ may receive angular momentum corrections.
    \needscheck

  \item \textbf{Sen formula graviton contribution.}
    The $-424$ in Sen's formula includes subtleties with the
    graviton zero modes on the Euclidean instanton.  In SCT, the
    graviton propagator is modified (the $\Pi_{\mathrm{TT}}$ pole
    shifts), which may alter the zero-mode counting.
    \gap

  \item \textbf{Modified evaporation endpoint.}
    The late-stage evaporation (where $T_H \sim \Lambda$) requires
    the full nonlocal Regge--Wheeler equation~\eqref{eq:RW-SCT}
    and connects to MR-9.  \gap

  \item \textbf{Information paradox.}
    The spectral action's entire-function property may provide a
    natural resolution of the information paradox through modified
    near-horizon physics.  This is deferred to MR-9.  \gap
\end{enumerate}


%%======================================================================
\section{Master Reference Table}
\label{sec:refs}
%%======================================================================

\begin{center}
\begin{longtable}{p{1.2cm}p{5.2cm}p{3.5cm}p{4cm}}
\toprule
\textbf{Ref} & \textbf{Citation} & \textbf{Key result} & \textbf{Use in MT-1} \\
\midrule
\endhead
R1 & Wald (1993), Phys.\ Rev.\ D \textbf{48}, R3427
  & BH entropy = Noether charge
  & Wald formula~\eqref{eq:wald-general} \\
R2 & Iyer--Wald (1994), Phys.\ Rev.\ D \textbf{50}, 846
  & First law for diff-invariant theories
  & First law~\eqref{eq:first-law} \\
R3 & Jacobson--Myers (1993), PRL \textbf{70}, 3684
  & Entropy of Lovelock gravity
  & JM formula~\eqref{eq:JM-formula} \\
R4 & Jacobson--Kang--Myers (1994), gr-qc/9312023
  & $C^2$ and $R^2$ Wald entropy
  & $\delta S_{C^2}$, $\delta S_{R^2}$ \\
R5 & Sen (2012), arXiv:1205.0971
  & $c_{\log}$ for non-extremal BH
  & Eq.~\eqref{eq:sen-formula} \\
R6 & Wall (2015), arXiv:1504.08040
  & Second law for HD gravity
  & Section~\ref{sec:four-laws} \\
R7 & Bardeen--Carter--Hawking (1973)
  & Four laws of BH mechanics
  & Zeroth, first, third laws \\
R8 & Hawking (1975), CMP \textbf{43}, 199
  & Hawking radiation
  & Section~\ref{sec:hawking} \\
R9 & Regge--Wheeler (1957)
  & Perturbation equation
  & Greybody factors \\
R10 & Buoninfante--Mazumdar (2020), 1903.01542
  & Nonlocal gravity + BH thermo
  & Local limit on horizon \\
R11 & Conroy--Edholm--Mazumdar (2017), 1509.01636
  & Nonlocal Wald entropy
  & Nonlocal corrections \\
R12 & Christensen--Duff (1978), NPB \textbf{154}, 301
  & Conformal anomaly coefficients
  & Relation to $c_{\log}$ \\
R13 & Stelle (1977), PRD \textbf{16}, 953
  & Higher-derivative gravity
  & Local limit benchmark \\
R14 & Mann--Solodukhin (1998), hep-th/9709064
  & Conical singularity method
  & Alternative $c_{\log}$ derivation \\
R15 & Bombelli et al.\ (1986), PRD \textbf{34}, 373
  & Entanglement entropy = area law
  & Entanglement approach \\
R16 & Padmanabhan (2010), \textit{Gravitation}
  & Weyl tensor on horizons
  & Prop.~\ref{prop:weyl-binormal} \\
R17 & Kaul--Majumdar (2000), PRL \textbf{84}, 5255
  & LQG $c_{\log} = -3/2$
  & Comparison table \\
R18 & Connes--Chamseddine (2010), 1004.0464
  & Spectral action and gravity
  & Spectral framework \\
\bottomrule
\end{longtable}
\end{center}


%%======================================================================
\section{Notation Summary}
\label{sec:notation}
%%======================================================================

\begin{center}
\begin{tabular}{lll}
\toprule
Symbol & Definition & Value/Source \\
\midrule
$\alpha_C$ & Weyl${}^2$ coefficient & $13/120$ (NT-1b Phase~3) \\
$\alpha_R(\xi)$ & $R^2$ coefficient & $2(\xi - 1/6)^2$ \\
$c_{\log}$ & Logarithmic entropy correction & $37/24$ (Sen + SM, microcanonical) \\
$\kappa$ & Surface gravity & $1/(4GM)$ (Schwarzschild) \\
$T_H$ & Hawking temperature & $1/(8\pi GM)$ \\
$r_H$ & Horizon radius & $2GM$ \\
$A$ & Horizon area & $16\pi G^2 M^2$ \\
$\hat{\epsilon}_{\mu\nu}$ & Bifurcation surface binormal &
  $\hat{\epsilon}_{ab}\hat{\epsilon}^{ab} = -2$ \\
$\chi(\mathcal{H})$ & Euler characteristic & $2$ for $S^2$ \\
$N_s, N_f, N_v$ & SM field counts & $4, 45, 12$ \\
\bottomrule
\end{tabular}
\end{center}


%%======================================================================
\section{Summary: The Full BH Entropy Formula in SCT}
\label{sec:summary}
%%======================================================================

Collecting all contributions, the black hole entropy in SCT for a
Schwarzschild BH is:

\begin{tcolorbox}[colback=blue!5!white, colframe=blue!50!black,
  title={\textbf{MT-1 Main Formula (to be verified)}}]
\begin{equation}\label{eq:full-entropy}
  S = \underbrace{\frac{A}{4G}}_{\text{Bekenstein--Hawking}}
  + \underbrace{\frac{13}{120\pi}}_{\text{Wald ($C^2$, topological)}}
  + \underbrace{\frac{37}{24}\,
    \ln\!\left(\frac{A}{\ell_P^2}\right)}_{\text{one-loop (Sen)}}
  + \mathcal{O}(1).
\end{equation}
\end{tcolorbox}

\noindent
The three terms have distinct physical origins:
\begin{enumerate}[nosep]
  \item \textbf{$A/(4G)$}: classical area law from the
    Einstein--Hilbert sector of the spectral action (Wald entropy).
  \item \textbf{$13/(120\pi)$}: tree-level correction from the
    Weyl-squared sector; topological (depends on $\chi(\mathcal{H})$,
    not on~$A$); parameter-free via $\alpha_C = 13/120$.
  \item \textbf{$(37/24)\ln(A/\ell_P^2)$}: one-loop quantum
    correction from all SM fields + graviton; parameter-free via
    SM field content ($n_S = 4$, $n_F = 22.5$ Dirac, $n_V = 12$).
\end{enumerate}

For an astrophysical BH ($M = 10\,M_\odot$):
$S_{\mathrm{BH}} \approx 1.05 \times 10^{79}$,
$\delta S_{\mathrm{Wald}} \approx 0.035$,
$\delta S_{\mathrm{log}} \approx 283$.
The hierarchy is $S_{\mathrm{BH}} \gg \delta S_{\mathrm{log}} \gg
\delta S_{\mathrm{Wald}}$.

\vfill
\begin{center}
\small\textit{End of MT-1 Literature Extraction.
All formulas extracted from cited sources; no original derivations.}
\end{center}

\end{document}
